% =====================================================================
%  Chapter 17c — The forcing-floor class theorem, and where it stops (S60-S62)
% =====================================================================
\chapter{The Forcing-Floor Class Theorem, and Where It Stops}
\label{ch:forcing-floor}

\begin{record}
\textbf{Sources of record.} Two session logs and one scoping handover, read in
full:

\smallskip
{\raggedright\ttfamily
2026-09-13-S60-cd1-forcing-floor-general.md\\
HANDOVER-S61-CD2.md\\
2026-09-16-S62-cd2-fractional-and-second-instance.md\par}

\smallskip
\Sn{61} has no session-log file, by design: its commit was a scoping handover
for \Sn{62}, not a session that touched the proof document or ran anything.

\smallskip
Each is cross-checked here against its \texttt{PROGRESS.md} \textsc{key
finding} block (\Sn{60}, \Sn{62}), and, where finer technical detail was
needed, against the proof document directly: \S\,\texttt{sec:cd1-general-floor}
(\Sn{60}). \Sn{62} made no \texttt{proofs/} edit, so there is nothing further
to check there; its findings live only in its own session log and its
\texttt{PROGRESS.md} block.

\smallskip
\textbf{Citation convention, as everywhere in this book.} Objects of the proof
document are named as text --- \texttt{thm:cd1-general-floor} and the rest ---
because that document is not this one and its labels do not resolve here.
External papers are named by author and arXiv identifier in running prose,
never by a \texttt{\textbackslash cite} key this book's own bibliography does
not carry. Only cross-references \emph{inside} this book are set as
references.

\smallskip
\textbf{No numerics.} Like the arc of Chapter~\ref{ch:pressure-testing}, this
one produced no experiment and no \texttt{results.db} row. Everything below is
derivation and reading, and every quantitative statement traces to a session
log, to \texttt{PROGRESS.md}, or to a labelled object in the proof document.
\end{record}

Chapter~\ref{ch:pressure-testing} closed with Wall~1 exactly where \Sn{47}
left it, tested three further times and unmoved. A fourth test followed it
directly: \Sn{59}, a second strategic audit, examined the programme's own
founding thermodynamic thesis --- Route~1 with \(\mu(T)\), Route~2 with the
auxiliary scalar \(\theta\), the scalar-\(T\) strong maximum principle itself
--- and found that it inherits Wall~1 \emph{by construction}, because Wall~1
is a property of \(\operatorname{div}u=0\) and both routes enforce
\(\operatorname{div}u=0\) exactly by design. Its \texttt{PROGRESS.md} entry
calls this the programme's sixth independent confirmation that the obstruction
is real. This chapter is not about \Sn{59}, and does not re-argue it; it is
recorded here only because it is the last event before the one this chapter
is about.

\section{The pivot, stated first}
\label{sec:ff-thesis}

\begin{record}
\textbf{The chapter's thesis, stated as its own claim rather than left
implicit.} After four more independent confirmations that Wall~1 does not
move, the owner changed the question being actively pursued. Not ``prove (A)
or (B)'' --- that question, and Wall~1's classification, are exactly where
\Sn{47} left them, and nothing in this chapter touches either. A different,
bounded, well-posed question about somebody else's already-published
mathematics: \emph{how much forcing do existing forced blow-up constructions
actually require, and is any of it removable?} This is Fefferman's
alternatives (C) and (D), studied from outside, never by putting a forcing
term into this programme's own arguments --- a standing, absolute rule
(Wall~9, \S\ref{sec:wall9}) that nothing below relaxes.

\smallskip
Three sessions answer parts of that question, and together they are a
complete small arc. \Sn{60} generalises one profile's forcing floor
(\Sn{55}, \S\ref{sec:pt-s55}) into a genuine class theorem --- the first
piece of general mathematics this line produced, rather than a check against
a single external object. \Sn{62}, after a scoping-only \Sn{61}, pushes that
theorem to its two natural next tests: extend it to fractional dissipation,
and find a second instance to run it against. Both are pushed as far as they
honestly go within three sessions. The extension fails for a reason that can
be written down exactly rather than merely asserted; the one further
instance the literature contains satisfies the theorem's hypotheses but sits
in a corner where the theorem has nothing sharp to say. The line is closed in
consequence. That is this chapter's ending, not a setback to argue around,
and the chapter says so as plainly as Chapter~\ref{ch:pressure-testing} said
of its own five-session arc that nothing moved.
\end{record}

\section{\texorpdfstring{\Sn{60}: a class theorem, and what it does not need}{S60: a class theorem, and what it does not need}}
\label{sec:ff-s60}

\Sn{60} was commissioned with two threads. Thread~1 asked whether \Sn{55}'s
one-profile forcing floor generalises. Thread~2 asked, separately, whether the
unforced OpenAI Euler construction of Chapter~\ref{ch:openai} survives being
given a fixed viscosity \(\nu>0\). The two are logically independent and the
session kept them so; this section does the same, and reports Thread~2 as the
closed matter it is before Thread~1 as the matter this chapter continues.

\subsection{The theorem}
\label{sec:ff-s60-theorem}

For axisymmetric, smooth \((u,p)\) with \(\operatorname{div}u=0\) and
\(\nu\ge0\), define the residual
\(f:=\partial_tu+(u\cdot\nabla)u-\nu\Delta u+\nabla p\) --- literally what is
left over once \(u\) is substituted into the axisymmetric momentum equation,
Navier--Stokes at \(\nu>0\) or Euler at \(\nu=0\). With \(\Gamma=ru_\vartheta\),
\(M(t)=\max\Gamma>0\) attained on a set \(\Sigma(t)\) satisfying the uniform
compactness carried over from \Sn{55}'s \texttt{lem:wpt3-forced-max},
\(R(t)=\max_{\Sigma(t)}r\) and \(F(t)=\max_{\Sigma(t)}\abs{f_\vartheta}\),
\texttt{thm:cd1-general-floor} gives three equivalent-strength forms. The
first needs no geometric input of any kind:
\begin{align}
\text{(i)}\quad & \frac{f_\vartheta(x_*)}{u_\vartheta(x_*)}\;\ge\;
  \frac{\dee^+}{\dee t}\log M(t) \qquad\text{for some }x_*\in\Sigma(t), \\[4pt]
\text{(ii)}\quad & F(t)\;\ge\;\frac{1}{R(t)}\,\frac{\dee^+}{\dee t}M(t), \\[4pt]
\text{(iii)}\quad & \int_s^tR(\sigma)F(\sigma)\dee\sigma\;\ge\;M(t)-M(s)
  \qquad\text{(no differentiability of }M\text{ used anywhere).}
\end{align}
Form (i) is the one worth carrying forward. It has no core radius, no
prescribed growth rate and no self-similarity in it, and reads: \emph{at a
swirl maximum, the azimuthal residual measured against the local swirl is at
least the logarithmic growth rate of the swirl maximum.} That sentence is the
general form of what \Sn{55} could establish only for one profile, by
comparing four exponents of that profile's own azimuthal momentum equation.
\texttt{rem:cd1-euler} records that neither this theorem nor the \Sn{55}
lemmas it rests on use the sign of \(\nu\): the only place \(\nu\) appears in
the proof is the inequality \(\nu(\partial_r^2\Gamma+\partial_z^2\Gamma)\le0\)
at an interior maximum, which at \(\nu=0\) reads \(0\le0\). So the floor
covers forced axisymmetric Euler exactly as it covers forced Navier--Stokes
--- which matters for this line specifically, because more than one
construction the field is trying to de-force is Euler, not Navier--Stokes.

\subsection{The rate corollary, and the OpenAI recovery}
\label{sec:ff-s60-rate}

Adding two hypotheses turns the ratio form into a number. With (G1)
\(M(t)\ge c_1\tau^{-a}\), \(a>0\), a prescribed swirl growth rate, and (G2)
\(R(\sigma)\le R_0\tau^b\), a prescribed core-radius geometry, a bound
\(\abs{f_\vartheta}\le K\tau^{-p}\) at the maximisers is forced by
\texttt{cor:cd1-rate} to satisfy
\[
p\;\ge\;1+a+b, \qquad\text{and at equality}\qquad K\;\ge\;\frac{a\,c_1}{R_0}.
\]
\texttt{cor:cd1-recovers} checks this against the OpenAI leading axisymmetric
field, whose own swirl maximum and maximiser radius (\Sn{55}'s
\texttt{eq:wpt3-Mexact}, \texttt{eq:wpt3-rm}) are exactly (G1)--(G2) with
\(a=h\), \(c_1=c_*\), \(b=1/2\), \(R_0=c_r\): the corollary forces
\(p\ge3/2+h\) and \(K\ge hc_*/c_r\), which is \texttt{prop:wpt3-floor}'s own
exponent and constant, recovered exactly. The general theorem contains the
\Sn{55} result rather than merely resembling it.

\subsection{The correction to \texorpdfstring{\Sn{55}}{S55}'s own record}
\label{sec:ff-s60-correction}

\Sn{55}'s log said, of the exact power \(M(t)=c_*\tau^{-h}\), that this
exactness --- not \(\asymp\) --- was ``what makes the floor exact rather than
a scaling heuristic, and it is available only because the leading field is
exactly self-similar.'' \Sn{60} found that overstated, and \texttt{rem:cd1-%
selfsimilar} splits the claim three ways. The exponent needs none of it: the
bound \(p\ge1+a+b\) follows from the one-sided asymptotics (G1)--(G2) alone,
through the \emph{integrated} form (iii), which needs no differentiability of
\(M\) at any point. The constant needs none of it either: \(K\ge ac_1/R_0\) is
recovered from the asymptotic constants alone. What exactness actually buys
is time-localisation, and only that --- with it, the floor holds at every
\(t\), at an identified point; without it, one gets only ``for every \(s<T\)
there is some \(\sigma\in(s,T)\) at which the floor holds,'' with \(\sigma\)
unidentified. The one genuinely profile-specific step of \Sn{55} was deriving
the maximiser radius as an \emph{output} of the exact profile; in the general
theorem that radius is an \emph{input}, (G2), and form~(i) dispenses with it
entirely by measuring the residual against the local swirl rather than in
absolute units. This is a real correction to how the programme's own record
characterised its earlier result, in the same spirit as the corrections
Chapter~\ref{ch:pressure-testing} records elsewhere: right about the
conclusion, imprecise about the reason.

\subsection{The hypothesis audit, and the wall that does not move}
\label{sec:ff-s60-audit}

\texttt{rem:cd1-audit} sorts the theorem's hypotheses into three groups.
Smoothness and an orientation convention are generic and do no work. The
pressure-free swirl equation and the exact vanishing of the singular
\(-2\nu/r\) diffusion term at an interior maximum are also generic, in a
stronger sense: both follow from axisymmetry and \(\operatorname{div}u=0\)
alone, and need nothing else --- in particular the maximiser is automatically
off the axis, since \(\Gamma=O(r^2)\) there, so no \(\varepsilon\)-exhaustion
is needed to reach it. The one hypothesis that must be checked per
construction is (H2), uniform compactness of the maximiser set \(\Sigma(t)\):
it is what lets a Danskin argument extract the upper bound at all, it is
automatic for a compactly supported field, and it can fail if the swirl
maximum is allowed to escape to spatial infinity.

One boundary is not a matter of technique at all. \emph{Axisymmetry is the
only lever the theorem uses}, through \((\nabla p)_\vartheta=
r^{-1}\partial_\vartheta p=0\). Outside that class there is no scalar swirl
equation, and no analogue is proved or claimed. This is why
\texttt{thm:cd1-general-floor}, exactly like \texttt{prop:wpt3-floor} before
it, says nothing at all about the field OpenAI's own Theorem~1.1 actually
claims a blow-up for --- that field is non-axisymmetric. Wall~9 governs and
its classification does not move.

\subsection{A positive-measure floor, and where it cannot reach}
\label{sec:ff-s60-lq}

\Sn{55} recorded, and deliberately did not pursue, whether a uniform
(positive-measure) version of the floor is provable. \texttt{prop:cd1-Lq}
answers one specific version of that in four lines: for compactly supported
\(\Gamma\), testing the swirl equation with \(\abs{\Gamma}^{q-2}\Gamma\), the
drift integrates away by \(\operatorname{div}u=0\), the diffusion contributes
a sign-definite sink, the same singular term vanishes identically for the
same structural reason as \S\ref{sec:ff-s60-audit}, and H\"older finishes it:
\[
\frac{\dee}{\dee t}\norm{\Gamma(\cdot,t)}_{L^q}\;\le\;\norm{rf_\vartheta(\cdot,t)}_{L^q},
\qquad 2\le q<\infty,\ \nu\ge0.
\]
Three honest caveats travel with it. It bounds \(\norm{rf_\vartheta}_{L^q}\),
not \(\norm{f_\vartheta}_{L^q}\); converting the two needs a radius bound on
the \emph{support} of the residual, strictly stronger than (G2). It does not
contradict \texttt{rem:wpt3-pointwise}(b) --- a different hypothesis, growth
of a norm rather than of a pointwise maximum, and the spike counterexample
that motivated that remark still applies to the pointwise statement. And it
does \textbf{not} apply to the OpenAI leading field at all: there
\(\Gamma^{(0)}\to\sqrt2\,c_\infty2^hr^{-2h}\) at large radius, independently
of \(q\), so it lies in no \(L^{q'}(\mathbb R^3)\) for any finite \(q'\). The
positive-measure floor is a tool for compactly supported constructions, and
the one profile this line had in hand at \Sn{60} was not one.

\subsection{Thread 2, closed: the unforced Euler mechanism against viscosity}
\label{sec:ff-s60-thread2}

Separately from all of the above, \Sn{60} asked whether the OpenAI \emph{Euler}
construction --- an iterated-oscillation scheme in the C\'ordoba--Mart\'inez-%
Zoroa lineage, not a self-similar profile, whose own Theorem~1.1 claims a
smooth compactly supported datum blowing up in finite time for the unforced
Euler equations --- survives being given a fixed viscosity. Read from the
primary source, the packet construction is available only under the paper's
own condition (3.12), which combined with its own (5.10)/(5.12) forces the
packet frequency at stage \(j\) to exceed a fixed enormous power of the
gradient it produces: \(k_j>h_j^{10^8Q}\) for a fixed constant \(Q\). Viscous
marginality --- the packet's dissipation rate must not exceed the stretching
rate that amplifies it --- instead needs \(k_j\lesssim\ell_j(h_{j-1}/\nu)^{1/2}\).
The two are incompatible in the exponent, by a factor of order \(2\times10^8Q\),
for every fixed \(\nu>0\) once \(h_j\to\infty\); the computation is explicit
and was checked against the paper's own displayed scales rather than
estimated.

This structural negative is corroborated three independent ways, all from
primary sources fetched and read that session: the same construction lineage,
pushed by C\'ordoba, Mart\'inez-Zoroa and Zheng (arXiv:2407.06776) to
hypodissipation \(\abs{\nabla}^\alpha\), survives only to
\(\alpha<\alpha_0=(22-8\sqrt7)/9\approx0.093\) out of a real viscosity's
\(\alpha=2\) --- and still needs forcing; OpenAI's own companion
Navier--Stokes paper abandons this packet geometry entirely for a
self-similar background sitting exactly at the viscous scale
(\(\ell_r\asymp\tau^{1/2}\), \(\mathrm{Re}_r=O(1)\), viscous damping a leading
term throughout) when it needed viscosity to be survivable at all, and it too
needs forcing there; and the Euler paper's only appearance of \(\nu\) is an
auxiliary regularisation, introduced and then explicitly removed, that carries
no information about a physical viscous limit.

What this does \emph{not} do is answer the question that would actually
matter: whether the construction's smooth Euler datum, run through
Navier--Stokes at \(\nu\to0\), stays regular. That is exactly Constantin's own
named open problem, ``a gap between the two ranges of viscosities'' with ``no
[known way] to close it,'' and closing it would need a uniform-in-\(\nu\) a
priori estimate --- the same class of object Wall~1 obstructs. \Sn{60}'s own
recommendation was explicit: this is a complete survey, not a work package,
and opening it further would be manufacturing a target. Thread~2 is not
continued anywhere in this chapter, and no further session touched it.

\section{\texorpdfstring{\Sn{61}: scoping the next step}{S61: scoping the next step}}
\label{sec:ff-s61}

\Sn{61} is a handover, not a session: no worktree, no proof-document edit, no
session log, by design. It set two deliverables for \Sn{62}, in a deliberate
order. \emph{Do the harder, more likely to fail thing first}: attempt to
extend the swirl maximum principle underlying \texttt{thm:cd1-general-floor}
from the classical Laplacian to fractional dissipation, \(\abs{\nabla}^\alpha\)
with \(0<\alpha<2\), because if that fails there is no point searching for
constructions that would need it. Only then look for a second instance of the
theorem's hypotheses in the literature, checking any candidate from its
primary source rather than from a secondary description, exactly as \Sn{54}
and \Sn{55} had done for the OpenAI profile.

Five named failure modes carried forward from \Sn{60} governed both
deliverables, and are worth restating because \Sn{62} in fact triggered two of
them and treated both as legitimate outcomes rather than setbacks: \textbf{H1}
--- the class has one member, and no second instance exists to test the
theorem against; \textbf{H2} --- the fractional swirl equation has no clean
scalar form; \textbf{H3} --- rediscovering a result already in the fractional
Calder\'on--Zygmund literature without attribution, guarded against by
reading first and deriving second; \textbf{H4} --- drifting into advocacy for
(C)/(D) as a goal in its own right, which this programme does not do;
\textbf{H5} --- blurring the standing prohibition on introducing forcing into
this programme's own arguments, which nothing in \Sn{60}--\Sn{62} does.

\section{\texorpdfstring{\Sn{62}: the natural extension, and the second instance}{S62: the natural extension, and the second instance}}
\label{sec:ff-s62}

\subsection{Deliverable 1: the fractional extension fails, precisely}
\label{sec:ff-s62-fractional}

The classical (\(\alpha=2\)) proof works because the maximum principle is
applied not to \(u_\vartheta\) but to \(\Gamma=ru_\vartheta\), via the
vector-Laplacian identity \((\Delta u)_\vartheta=\Delta
u_\vartheta-u_\vartheta/r^2\), which converts the momentum equation for
\(u_\vartheta\) into a genuinely \emph{closed scalar equation for \(\Gamma\)}.
For a pure-swirl axisymmetric field, the azimuthal component of the fractional
Laplacian \((-\Delta)^{\alpha/2}u\) --- independently re-derived this session,
component-wise in the fixed Cartesian frame --- works out instead to
\[
g_\vartheta(x)\;=\;c_{3,\alpha}\,\mathrm{P.V.}\!\int
\bigl[u_\vartheta(x)-u_\vartheta(y)\cos(\theta_x-\theta_y)\bigr]\,
\abs{x-y}^{-3-\alpha}\dee y.
\]
This is genuinely different from the classical Calder\'on-type scalar
comparison \(\theta(x)-\theta(y)\) that the C\'ordoba--C\'ordoba inequality
governs --- the PNAS~2003 result was fetched and confirmed to apply to
\emph{scalars} specifically --- because of the angle-dependent,
sign-changing weight \(\cos(\theta_x-\theta_y)\).

A concrete counterexample shows the naive extension fails outright, not
merely that a proof was not found. Fix \(x\) at an interior maximiser of
\(\Gamma\), and take \(y\) at the \emph{same} angle (\(\cos(\theta_x-\theta_y)
=1\)) with radius \(r_y\) slightly less than \(r_x\). Since \(\Gamma'(r_x)=0\)
at the maximum, \(\Gamma(y)\approx\Gamma(x)\) to leading order, but
\(u_\vartheta(y)=\Gamma(y)/r_y>\Gamma(x)/r_x=u_\vartheta(x)\) purely because
\(r_y<r_x\) --- so the bracket \(u_\vartheta(x)-u_\vartheta(y)\) is
\emph{strictly negative}, even at zero angular separation, at a point where
the classical argument needs it to vanish or help. The mechanism that makes
the classical case work --- the local, \(y\to x\) degeneration of the angular
weight into a single algebraic curvature term that cancels exactly at
\(\Gamma\)'s critical point --- has no analogue once \(\alpha<2\) makes the
operator genuinely long-range: \(g_\vartheta\) cannot be re-expressed as a
nonlocal operator on \(\Gamma\) alone carrying a sign property at \(\Gamma\)'s
maximum. This is H2, precisely located rather than merely invoked: the
classical proof's specific mechanism --- reduction to a closed scalar
equation on \(\Gamma\) via the vector-Laplacian identity --- has no fractional
counterpart, because \(u_\vartheta\) and \(\Gamma\) need not share a
maximiser once the comparison point \(y\) is allowed to range freely in
radius.

Two things this is not. It is not a claim that no fractional swirl maximum
principle exists by any means at all --- an added radial-monotonicity
hypothesis on \(u_\vartheta\) near the maximiser would block exactly this
counterexample, and whether some other route closes the gap is a separate,
unattempted question. And it is not a rediscovery left unattributed: the
literature search run that session found no prior result addressing this
specific question. The consequence for this line is direct: no fractional
\texttt{cor:cd1-rate} exists, so nothing built earlier in this chapter carries
over to hypodissipative constructions.

\subsection{Deliverable 2: two ruled out, one found in a degenerate corner}
\label{sec:ff-s62-second}

Both candidates \Sn{60} had named were checked from primary sources and both
were ruled out --- on axisymmetry specifically, not on existence, and both
papers are confirmed real and were previously used correctly in this
programme for other purposes. One further instance, not named in the
handover, was found by a genuine search and satisfies the theorem's
hypotheses; it is degenerate in the one respect that matters for a numeric
floor. Table~\ref{tab:ff-candidates} records all three.

\begin{table}[ht]
\centering
\footnotesize
\begin{tabular}{@{}>{\raggedright\arraybackslash}p{0.19\textwidth}
                     >{\raggedright\arraybackslash}p{0.23\textwidth}
                     >{\raggedright\arraybackslash}p{0.23\textwidth}
                     >{\raggedright\arraybackslash}p{0.21\textwidth}@{}}
\toprule
Paper & What it is & Axisymmetric? & Applies to the floor theorem? \\
\midrule
arXiv:2309.08495\newline(C\'ordoba--%
  Mart\'inez-Zoroa)
  & Forced 3D Euler, \(C^{1,1/2-\varepsilon}\cap L^2\) force
  & \textbf{No} --- its own abstract states ``a specific focus on
    non-axisymmetric solutions''
  & Excluded \\[8pt]
arXiv:2407.06776\newline(C\'ordoba--%
  Mart\'inez-Zoroa--Zheng)
  & Forced hypodissipative NS, \(\abs{\nabla}^\alpha\),
    \(\alpha<\alpha_0\approx0.093\)
  & \textbf{No} --- \S1.2.5 states verbatim the vorticity is
    \emph{axial} (reflection symmetry across coordinate planes),
    velocity \emph{polar}: a Cartesian vortex-layer symmetry, not
    azimuthal swirl
  & Excluded --- completes, not corrects, \Sn{60}, which used this
    paper only for its Thread-2 threshold \\[8pt]
arXiv:2311.12306\newline(Qi S.\ Zhang)
  & Forced axisymmetric pure-swirl ASNS, two explicit closed-form
    self-similar solutions (Theorem~1.1)
  & \textbf{Yes} --- its system (1.4) is literally the ASNS system
    \texttt{lem:wpt3-forced-swirl} is built for; a finite-cylinder
    domain makes (H2) automatic
  & Hypotheses satisfied; degenerate \(a=0\) corner of
    \texttt{cor:cd1-rate} (Table~\ref{tab:ff-rate}) \\
\bottomrule
\end{tabular}
\caption[The three candidates checked at S62]{The two constructions named in
\texttt{HANDOVER-S61-CD2.md}, both real and both previously used correctly in
this programme, ruled out here on a dimension neither prior use had checked;
and the one further instance a genuine search turned up.}
\label{tab:ff-candidates}
\end{table}

The new instance repays derivation. Zhang's self-similar solution is
\(v_\vartheta(r,t)=u_\vartheta(r,t)+\alpha r\), with
\(u_\vartheta(r,t)=\varphi_0(\rho)/\sqrt{2(T-t)}\) and
\(\rho=r/\sqrt{2(T-t)}\) (the paper's own equations). Substituting
\(r=\rho\sqrt{2(T-t)}\),
\[
\Gamma(r,t)\;=\;r\,v_\vartheta\;=\;\rho\,\varphi_0(\rho)\;+\;2\alpha(T-t)\rho^2 .
\]
The first term is an exact function of \(\rho\) alone; the second is a
genuine, but asymptotically subleading, \(O(T-t)\) correction, vanishing at
any fixed \(\rho\) as \(t\to T\) --- a real residual, not the ``no residual
\(t\)-dependence'' the session's own subagent draft first claimed, a
distinction the orchestrator's independent re-derivation flagged and
corrected in the write-up without changing the substance. Because
\(\varphi_0(r)=O(r/(r^2+1))\) decays (the paper's own bound),
\(\rho\varphi_0(\rho)\) attains a finite maximum at some \(\rho_*=O(1)\), and
its physical location \(r_m(t)=\rho_*\sqrt{2(T-t)}\to0\) --- exponent
\(b=1/2\) in this document's (G1)/(G2) notation, matching OpenAI's own \(b\)
exactly. At that maximiser the correction term is \(O(T-t)\to0\), so
\(M(t)=\max\Gamma\to\rho_*\varphi_0(\rho_*)\), a nonzero \emph{constant} ---
exponent \(a=0\).

\begin{table}[ht]
\centering
\small
\begin{tabular}{@{}>{\raggedright\arraybackslash}p{0.30\textwidth} c c
                     >{\raggedright\arraybackslash}p{0.36\textwidth}@{}}
\toprule
Construction & \(a\) & \(b\) &
  Floor from \texttt{cor:cd1-rate} \\
\midrule
OpenAI leading axisymmetric field (\Sn{55}, \Sn{60})
  & \(h\) (\(<1/100\)) & \(1/2\)
  & \(p\ge\tfrac32+h\), \(K\ge hc_*/c_r\) --- nontrivial, matches the
    source's own exponent exactly \\[6pt]
Zhang ASNS self-similar solution (\Sn{62})
  & \(0\) & \(1/2\)
  & \(p\ge\tfrac32\), \(K\ge0\) --- the trivial floor \\
\bottomrule
\end{tabular}
\caption[The rate corollary, applied twice]{\texttt{cor:cd1-rate} run against
both known instances. The class of fields satisfying
\texttt{thm:cd1-general-floor}'s hypotheses has (at least) two members; the
class for which the rate corollary says something \emph{nontrivial}
(\(a>0\)) still has exactly one.}
\label{tab:ff-rate}
\end{table}

Because \(a=0\), \texttt{cor:cd1-rate} yields only \(K\ge0\) here --- a floor
every force trivially satisfies. Zhang's construction satisfies the
theorem's hypotheses and is a genuine second member of the class the theorem
speaks about; it is not a second member of the smaller class for which the
theorem says anything a reader would call sharp. This is the precise,
corrected form of failure mode H1: not the flat ``the class has one member''
named in advance as the most likely outcome, but a class with at least two
confirmed members split by whether the rate corollary is informative on
them. There is nothing left to compare for Zhang's case in the three-way
sense \Sn{55} ran for OpenAI's: a trivial floor has no source-stated forcing
scaling to check it against.

\section{What the arc changed, and where it stopped}
\label{sec:ff-ledger}

\begin{figure}[ht]
\centering
\begin{tikzpicture}[
  >={Latex[length=2mm]},
  ev/.style={draw, rounded corners=2pt, font=\scriptsize, align=left,
             inner sep=4pt, text width=68mm},
  posbox/.style={ev, draw=provedgreen, thick},
  neg/.style={ev, draw=impossiblered, thick},
  neu/.style={ev, draw=rulegrey},
  ar/.style={->, thick, rulegrey}
]
\node[neu] (pivot) at (0,0)
  {\textbf{Pivot} (after \Sn{52}, \Sn{53}, \Sn{57}, \Sn{59}).
   Active question changes from ``prove (A)/(B)'' to ``how much forcing do
   existing forced constructions need, and is it removable?'' --- never by
   forcing this programme's own arguments};
\node[posbox, below=6mm of pivot] (s60a)
  {\Sn{60}, Thread 1. \texttt{thm:cd1-general-floor}: a genuine class
   theorem, \(\nu\ge0\) including Euler, recovering \texttt{prop:wpt3-floor}
   exactly as one instance};
\node[neu, below=6mm of s60a] (s60b)
  {\Sn{60}, Thread 2. OpenAI's unforced Euler mechanism structurally
   incompatible with any fixed \(\nu>0\) --- corroborated three ways,
   closed as a survey, not opened as a work package};
\node[neu, below=6mm of s60b] (s61)
  {\Sn{61} (scoping only). Attempt the fractional extension first; then
   search for a second instance};
\node[neg, below=6mm of s61] (s62a)
  {\Sn{62}, Deliverable 1. Fractional extension \textbf{fails} (H2),
   precisely located: no vector-Laplacian identity survives for
   \(\abs{\nabla}^\alpha\), \(\alpha<2\)};
\node[neu, below=6mm of s62a] (s62b)
  {\Sn{62}, Deliverable 2. Two named candidates excluded on axisymmetry;
   one new instance found (Zhang), degenerate (\(a=0\))};
\node[draw=impossiblered, very thick, rounded corners=2pt, font=\scriptsize,
      align=center, inner sep=5pt, text width=68mm, below=6mm of s62b]
  (closed) {\textbf{Line closed.} Recommendation taken; no further CD
   session};

\foreach \a/\b in {pivot/s60a, s60a/s60b, s60b/s61, s61/s62a, s62a/s62b, s62b/closed}
  {\draw[ar] (\a) -- (\b);}
\end{tikzpicture}
\caption[The three-session arc]{Pivot, general theorem, natural extension
attempted and priced, closed --- inside three sessions. The chain is drawn
straight because nothing in it doubled back: each session answered the
question the previous one posed and handed forward exactly one next
question.}
\label{fig:ff-arc}
\end{figure}

\begin{table}[ht]
\centering
\small
\begin{tabular}{@{}l p{0.34\textwidth} p{0.34\textwidth}@{}}
\toprule
Session & What it settled & What it left standing \\
\midrule
\Sn{60}
  & \texttt{prop:wpt3-floor} generalises into a genuine class theorem for
    \(\nu\ge0\), needing no geometric hypothesis in its sharpest form; the
    OpenAI Euler mechanism fails at fixed \(\nu>0\), structurally
  & Axisymmetry is the theorem's absolute boundary; the
    \(\nu\to0\) question is Constantin's own open gap, not tractable here \\[3pt]
\Sn{61}
  & Nothing new (scoping only)
  & Ordered the two live questions and named the failure modes both
    triggered in the event \\[3pt]
\Sn{62}
  & The fractional extension fails for an exactly locatable reason; a
    second axisymmetric forced instance exists and was checked
  & No fractional floor exists to build on; the nontrivial-floor class
    still has one member \\
\bottomrule
\end{tabular}
\caption[The arc's ledger]{Three sessions, one genuine theorem, one
precisely-located negative, one closed line.}
\label{tab:ff-ledger}
\end{table}

\begin{obsv}[A closed line is not a failed one]
\label{obs:ff-closed}
\Sn{62}'s own recommendation is unambiguous: close the line, and do not
manufacture a reason to continue when the honest answer is that it is done.
The owner accepted that recommendation. Nothing here is a setback dressed up
as a stopping point. \texttt{thm:cd1-general-floor} is real mathematics ---
elementary, by its own account, but general in a way nothing in this
programme's forcing-floor work was before it, and it survives being pushed:
it extends to \(\nu=0\) for free, it recovers its own special case exactly,
and a second, entirely independent construction (Zhang's) satisfies its
hypotheses on the first genuine attempt to find one. The fractional question
was a real question with a real answer, not a convenient place to stop
looking, and the answer is precisely locatable rather than a shrug: the
proof's one load-bearing algebraic fact --- that the vector-Laplacian
identity converts a vector equation for \(u_\vartheta\) into a scalar one for
\(\Gamma\) --- is a fact about the classical Laplacian specifically and has
no long-range analogue.

\smallskip
Nothing here is a step toward alternatives (A) or (B), and nothing here
argues that (C) or (D) ought to be anyone's goal. The programme's own
standing rule is unmoved by any of it: no forcing term has entered, or will
enter, an argument of this document toward its own object. What three
sessions produced instead is exactly what the pivot asked for --- a
sharper, numerically checkable answer to a question about other people's
already-published constructions --- and an honest report that the question,
for this theorem, has been answered as far as it currently can be.
\end{obsv}
